% ===================================================================
%  اللوحة رقم ٧ — القرية في الشتاء. — المزرعة في الربيع.
%  Traduction 2026 d'apres texto/io, controlee sur texto/fr, texto/en
%  et texto/es. Consigne : voir texto/ar/10-lawha-01.tex.
% ===================================================================

%%K t07-noto-1 noto
\VUnotes{143.2mm}{%
(*) لتفاصيل الوجبة انظر اللوحتين ٦ و١٣.%
}

%%K t07-apar-1 apar
\VUsaut{4.0mm}
\VUcentre{13.2pt}{90}{السلسلة الثانية}
\VUsaut{3.0mm}
\VUfilet{16mm}
\VUsaut{5.6mm}
\VUtitre{15.0pt}{60}{اللوحة رقم ٧}
\VUsaut{3.0mm}

%%K t07-tit-1 sub
\VUcentre{12.0pt}{0}{\textit{المشهد الأول.}}
\VUsaut{3.0mm}

%%K t07-tit-2 sub
\VUpk{10.2pt}{40}{القرية في الشتاء.}
\VUsaut{4.0mm}

%%K t07-c1-01-1 p
1. --- في عطلة رأس السنة ذهبت مع أبي إلى «القرية الجديدة»، وهي قرية
صغيرة كانت له فيها أعمال يقضيها.

%%K t07-c1-02-1 p
2. --- ركبنا \VUgras{عربة النقل} \textsuperscript{(11)} التي تسير بين المدينة
وتلك القرية؛ لكن البرد كان شديدًا، فلم تكن الرحلة في عربة بهذا القدر من
عدم الراحة سارّة أبدًا.

%%K t07-c1-03-1 p
3. --- ولذلك فرحنا فرحًا صادقًا حين رأينا \VUgras{السائق} يوقف عربته في
الساحة، أمام \VUgras{النُّزل} \textsuperscript{(2)} المسمى \textit{الدب الأبيض}.
وكان \VUgras{صاحب النُّزل} \textsuperscript{(4)} ينتظرنا على عتبته يدخّن
\VUgras{غليونه} في هدوء.

%%K t07-c1-03-2 p
دعانا إلى الدخول، ولم أحتج إلى أن يُدعى مرتين حتى استقررت أمام نار عيدان
الكرم المتطقطقة في الموقد. عندئذ لم أعبأ بـ\VUgras{الريح الشمالية}
اللاذعة التي كانت تُصرّ \VUgras{اللافتة} \textsuperscript{(3)} القديمة وتحمل في
دوامات سوداء \VUgras{الدخان} \textsuperscript{(1)} الخارج من المِدخنة.

%%K t07-c1-04-1 p
4. --- وكان وقت الغداء. فمن غير مزيد انتظار وفّينا حقّها للوجبة البسيطة
اللذيذة التي قُدّمت لنا (*).

%%K t07-tit-3 sub
\VUpk{10.2pt}{40}{بضع زيارات.}
\VUsaut{3.0mm}

%%K t07-c1-05-1 p
5. --- بعد الغداء ذهبت مع أبي، وهو وكيل تأمين، لزيارة زبائنه. زرنا أولًا
\VUgras{الحدّاد} \textsuperscript{(5)} الساكن إلى جانب النُّزل. كان مشغولًا
بنعل حصان. وقد أعجبتني قوة مساعده الذي كان يمسك \VUgras{حافر} الحصان
بإحكام على مئزره الجلدي بينما يثبّت الحدّاد \VUgras{النعل} \textsuperscript{(6)}
بضربات مطرقة ثقيلة.

%%K t07-c1-06-1 p
6. --- ثم ذهبنا إلى \VUgras{إسكافي} \textsuperscript{(7)} القرية، يعقوب العجوز.
ومع أن لافتته \VUgras{جزمة} رائعة، فقد كان يكتفي بتنعيل زوج قديم من أحذية
بالية.

%%K t07-c1-06-2 p
قال لنا الشيخ ضاحكًا: «في المدينة كنت ”إسكافيًا“ ولم يكن لي زبائن. وهنا
أنا ”مُرقِّع“ والعمل كثير».

%%K t07-c1-07-1 p
7. --- وحدّثنا عن جاره \VUgras{الحلّاق} \textsuperscript{(8)} الذي زبائنه كلهم
غير راضين. \VUgras{أمواسه} مثلومة دائمًا و\VUgras{مقصّاته} لا تقصّ كما
ينبغي أبدًا. لكنه الحلّاق الوحيد في القرية، فليس للناس بُدّ، بطبيعة
الحال، من أن يحلق لهم ويقصّ شعورهم. وهو فوق ذلك رجل بخيل غير لطيف.

%%K t07-c1-08-1 p
8. --- ولحسن الحظ ليس \VUgras{الجيران} الآخرون مثله. فـ\VUgras{صانع
العجلات} \textsuperscript{(9)}، مثلًا، رجل طيب مجتهد خدوم. ومن دواعي السرور أن
تعطيه عربة ليصلحها. فعمله متقن دائمًا.

%%K t07-c1-09-1 p
9. --- ولم يكن ليعقوب العجوز ما يشكوه من \VUgras{السرّاج} \textsuperscript{(10)}
أيضًا، وهو يعينه أحيانًا على خياطة \VUgras{الأعنّة} حين يشتد العمل.

%%K t07-c1-09-2 p
وسأله أبي عن أسرته. «الجميع بخير. حفيدتي في \VUgras{المدرسة}
\textsuperscript{(12)}. و\VUgras{معلمتها} \textsuperscript{(13)} راضية عنها جدًا؛ فهي
\VUgras{تلميذة} \textsuperscript{(14)} مجتهدة. ليست كذلك الشقيّ ابني؛ لا يفكر
إلا في اللعب. و\VUgras{المعلم} \textsuperscript{(15)} لا يستطيع أن يجعله يتعلم
شيئًا».

%%K t07-tit-4 sub
\VUsaut{3.0mm}
\VUpk{10.2pt}{40}{أحاديث القرية.}
\VUsaut{3.0mm}

%%K t07-c1-10-1 p
10. --- ثم حدّثنا الشيخ بأخبار القرية. فقد كانوا سيبنون مدرسة جديدة، لأن
التي في \VUgras{دار البلدية} صارت ضيقة أكثر مما ينبغي. وذلك هيّن، إذ يكفي
شراء جزء من حديقة \VUgras{الطحّان}. وفي ذلك خير كثير للرجل المسكين. فمنذ
بدأ البرد لم تستطع \VUgras{أحجار} \VUgras{الطاحونة} \textsuperscript{(16)} أن
تطحن القمح، إذ كان \VUgras{الدولاب} \textsuperscript{(17)} محبوسًا في
\VUgras{جليد} \textsuperscript{(25)} \VUgras{الجدول} \textsuperscript{(23)}.

%%K t07-c1-11-1 p
11. --- «وبمناسبة البناء، أرأيت \VUgras{كنيستنا} \textsuperscript{(18)} الجديدة؟
إنها بديعة المنظر تحت رداء الثلج. و\VUgras{الغربان} \textsuperscript{(19)} التي
لم تعد تعرف برجها القديم تحطّ في حزن على الشجرة الكبيرة في
\VUgras{المقبرة} \textsuperscript{(20)}».

%%K t07-c1-12-1 p
12. --- وذكْرُ المقبرة ذكّر الإسكافي بمن مات في السنة الماضية ممن كان أبي
يعرفهم. «كم من \VUgras{قبر} \textsuperscript{(21)}، يا سيدي العزيز، أُضيف إلى
سواه تحت \VUgras{السرو} \textsuperscript{(22)}! لقد أخذ الموت موثّقنا العجوز.
وكانت القرية كلها في \VUgras{جنازته}».

%%K t07-c1-13-1 p
13. --- «وها هو خليفته يتزحلق على الجدول. جاءني قبل قليل لأصلح له سير
\VUgras{زلاجته} \textsuperscript{(26)} الذي انقطع».

%%K t07-c1-14-1 p
14. --- «وهناك، على ضفة الجدول، \VUgras{حارس القرية} \textsuperscript{(27)}
الجديد. كان دركيًا، وهو رعب أشقياء القرية. يفرّون بمجرد أن يلمحوا
\VUgras{طماقه} \textsuperscript{(29)} و\VUgras{قبعته} \textsuperscript{(28)}».

%%K t07-c1-15-1 p
15. --- «آه! ها هو يوسف العجوز يجلب \VUgras{حِمل حطب} \textsuperscript{(30)}
للخبّاز. --- صدقت، قال أبي، أعرف رتل \VUgras{بغاله} \textsuperscript{(31)}. ولي
معه كلمة، وسأغتنم الفرصة. وداعًا يا يعقوب العجوز».

%%K t07-c1-16-1 p
16. --- وبينما كنا نخرج، كان أطفال القرية يغادرون المدرسة ويندفعون إلى
\VUgras{مزلقة الجليد} \textsuperscript{(33)}، مخاطرين بأن يقعوا ويكسروا عضوًا في
\VUgras{الوقعة} \textsuperscript{(34)}. وأحضر أحدهم \VUgras{مزلجته}
\textsuperscript{(32)} من عند صانع العجلات، وأجلس فيها صاحبًا له، وانطلق يعدو.

%%K t07-c1-17-1 p
17. --- واندفع آخرون إلى \VUgras{كومة ثلج} \textsuperscript{(37)} قرب
\VUgras{الجسر} \textsuperscript{(40)} وأخذوا يقذفون \VUgras{رجل الثلج}
\textsuperscript{(39)} الذي صنعوه في اليوم السابق، وقد ألبسوه قبعة وغليونًا
وحقيبة مدرسية على كتفه.

%%K t07-c1-18-1 p
18. --- ولم يكونوا يعبأون بالريح المتجمدة ولا بالبرد. فأكثرهم لم يكن له
حتى \VUgras{لِفاع} \textsuperscript{(35)}، ولكي يتدفأوا بعد معركة \VUgras{كرات
الثلج} \textsuperscript{(38)} كانت تكفيهم حفنة من \VUgras{الكستناء}
\textsuperscript{(42)} الساخنة يشترونها من جاكلين العجوز \VUgras{البائعة}، وهي
ترتجف من البرد تحت \VUgras{شالها} \textsuperscript{(43)} الرقيق أكثر مما ينبغي.

%%K t07-tit-5 sub
\VUsaut{3.0mm}
\VUcentre{12.0pt}{0}{\textit{المشهد الثاني.}}
\VUsaut{3.0mm}

%%K t07-tit-6 sub
\VUpk{10.2pt}{40}{المزرعة في الربيع.}
\VUsaut{4.0mm}

%%K t07-c2-01-1 p
1. --- على رغم ملذات الشتاء كلها، فإننا نستقبل عودة \VUgras{الربيع} بفرح
عميق.

%%K t07-c2-01-2 p
وما أبهج الصورة التي يصنعها فناء المزرعة في المساء، في شهر مايو!

%%K t07-c2-02-1 p
2. --- في الأفق تغيب \VUgras{الشمس} \textsuperscript{(1)} رويدًا، وقد أغرقت
\VUgras{الريف} \textsuperscript{(3)} بـ\VUgras{أشعتها} \textsuperscript{(2)} القرمزية.
و\VUgras{الحرّاث} \textsuperscript{(6)} المجدّ يسرع ليحرث \VUgras{حقله}
\textsuperscript{(4)}.

%%K t07-c2-02-2 p
وبـ\VUgras{محراثه} \textsuperscript{(5)} الذي يجرّه \VUgras{حصان} \textsuperscript{(7)}
متين، يخطّ \VUgras{الأخدود} \textsuperscript{(8)} الذي سينثر فيه \VUgras{البذّار}
\textsuperscript{(9)} بيديه \VUgras{البذر} الذي يعطي السنابل الذهبية.

%%K t07-c2-03-1 p
3. --- وقد قطع \VUgras{الحصّادون} \textsuperscript{(11)} بمناجلهم عشب المرج،
وجفّف الجامعون \VUgras{التبن} \textsuperscript{(10)} بتقليبه بـ\VUgras{المذاري}
\textsuperscript{(12)} والمشاطّ، وها هم عائدون.

%%K t07-c2-03-2 p
بعضهم يمشي أمام العربة الثقيلة التي تجرّها الثيران حاملين أدواتهم على
أكتافهم. وآخرون، وهم أكسل، استقروا في راحة على التبن الطيّب الرائحة.

%%K t07-c2-04-1 p
4. --- وهم في مرورهم يحيّون في ودّ \VUgras{البستاني} \textsuperscript{(16)}
المنشغل في \VUgras{البستان} \textsuperscript{(13)} بتقليم \VUgras{أشجار الفاكهة}
وتطعيم \VUgras{أشجار الكمثرى} \textsuperscript{(14)}. و\VUgras{أشجار الخوخ}
\textsuperscript{(15)} مزهرة، وفي \VUgras{حديقة الخضر} \textsuperscript{(17)} المسمّدة
جيدًا صارت الخضر والملفوف والخس في حال طيبة.

%%K t07-c2-04-2 p
وعلى الجدار المنخفض ينشر \VUgras{طاووس} \textsuperscript{(18)} جميل ذيله ليعرض
ريشه البديع.

%%K t07-tit-7 sub
\VUpk{10.2pt}{40}{فناء المزرعة.}
\VUsaut{3.0mm}

%%K t07-c2-05-1 p
5. --- أبواب \VUgras{الإسطبل} \textsuperscript{(19)} مفتوحة، ويُرى المِعلف
و\VUgras{الفراش} \textsuperscript{(23)} الخاص بـ\VUgras{الفرس} \textsuperscript{(21)}
التي فكّ عنها \VUgras{السائس} \textsuperscript{(20)} أدواتها لتوّه. و\VUgras{المُهر}
\textsuperscript{(22)}، وما أطرفه على ساقيه الطويلتين، يتبع أمه مرحًا.

%%K t07-c2-06-1 p
6. --- وقرب \VUgras{البئر} \textsuperscript{(29)} تذهب \VUgras{الأبقار}
\textsuperscript{(25)} لتشرب من \VUgras{الحوض} \textsuperscript{(26)} الخشبي الذي
يقوم لها مقام المشرب. و\VUgras{العِجل} \textsuperscript{(24)} يغتنم الفرصة
ليرضع، و\VUgras{الثور} \textsuperscript{(27)}، وقد شمّ رائحة العلف الطيبة، يخور
مبتهجًا ويرفع في زهو رأسه ذا \VUgras{القرنين} \textsuperscript{(28)} القويين.

%%K t07-c2-07-1 p
7. --- الجميع في العمل. \VUgras{الخادمة} \textsuperscript{(32)} تمسك \VUgras{حبل}
\textsuperscript{(31)} البئر. تسحب الماء في \VUgras{دلو} \textsuperscript{(30)} لتسقي
البهائم. وقرب \VUgras{الحظيرة} \textsuperscript{(33)} رجل يجمع التبن المتناثر،
وأمام باب \VUgras{بيت المزرعة} \textsuperscript{(34)} تغزل \VUgras{غازلة}
\textsuperscript{(35)} عجوز، بدولابها و\VUgras{مغزلها}، \VUgras{الكتّان} أو
\VUgras{القنّب} كما في الأيام الخوالي.

%%K t07-tit-8 sub
\VUsaut{3.0mm}
\VUpk{10.2pt}{40}{صاحب المزرعة.}
\VUsaut{3.0mm}

%%K t07-c2-08-1 p
8. --- \VUgras{صاحب المزرعة} \textsuperscript{(36)} يدير هذا العمل كله ويعطي
رجاله أوامرهم. يأمرهم بأن يضعوا في المأمن \VUgras{المِمشطة} \textsuperscript{(40)}
و\VUgras{المِعول} \textsuperscript{(39)} المتروكين قرب \VUgras{بيت} \textsuperscript{(38)}
\VUgras{الكلب} \textsuperscript{(37)}.

%%K t07-c2-09-1 p
9. --- وصياحه يفزع \VUgras{حمامة} \textsuperscript{(41)} تطير بأقصى سرعتها إلى
\VUgras{البُرج} \textsuperscript{(42)}، و\VUgras{الدجاجات} \textsuperscript{(43)} التي
تسرع عائدة إلى \VUgras{قنّ الدجاج} \textsuperscript{(44)}.

%%K t07-c2-10-1 p
10. --- وأمام \VUgras{حظيرة الغنم} \textsuperscript{(48)} ها هو \VUgras{الراعي}
\textsuperscript{(45)} العجوز يعيد \VUgras{قطيعه} \textsuperscript{(49)}. وهو يمسك
\VUgras{عصاه} \textsuperscript{(46)} التي لا تفارقه كما لا تفارقه
\VUgras{مِرْيَلته} \textsuperscript{(47)} الباهتة، يسوق إلى الحظيرة
\VUgras{الكباش} \textsuperscript{(50)} و\VUgras{النعاج} \textsuperscript{(53)}
و\VUgras{الحملان} \textsuperscript{(54)} و\VUgras{الماعز} \textsuperscript{(51)}
و\VUgras{الجداء} \textsuperscript{(52)}. و\VUgras{كلبه} \textsuperscript{(55)} الوفيّ
أرغوس يأتي في الآخر ويحثّ المتخلفين بنباحه.

%%K t07-c2-11-1 p
11. --- وهذا الضجيج والحركة كلها لا تعكّر هدوء \VUgras{البقرة الحلوب}
\textsuperscript{(56)} الطيبة، التي ترنّ \VUgras{جرسها} \textsuperscript{(57)} بينما
تُحلَب أمام باب \VUgras{زريبة البقر} \textsuperscript{(58)}.

%%K t07-c2-12-1 p
12. --- ويبدو أنها تفهم أن عليها أن تعطي لبنها ليتسنى لـ\VUgras{صاحبة
المزرعة} \textsuperscript{(60)} أن تصنع \VUgras{الزبدة} في \VUgras{المِمخضة}
\textsuperscript{(61)}، أو \VUgras{الأجبان} \textsuperscript{(64)} الممتازة التي
تُصفّى على اللوح الموضوع فوق \VUgras{الجفنة} \textsuperscript{(63)}.

%%K t07-c2-13-1 p
13. --- و\VUgras{مَلبنة} \textsuperscript{(59)} المزرعة كلها يحفظها نظيفة جدًا
\VUgras{أجير المزرعة} \textsuperscript{(65)} الذي غسل لتوّه \VUgras{أواني اللبن}
\textsuperscript{(62)} في الدلو الكبير الذي يحمله.

%%K t07-tit-9 sub
\VUsaut{3.0mm}
\VUpk{10.2pt}{40}{فناء الدواجن.}
\VUsaut{3.0mm}

%%K t07-c2-14-1 p
14. --- وهو في قباقيبه الثقيلة يمشي مشية خرقاء بعض الشيء، فيُطير
\VUgras{الديك الرومي} \textsuperscript{(68)} و\VUgras{البطات} \textsuperscript{(66)}
و\VUgras{ذكور البط} \textsuperscript{(67)} و\VUgras{الإوزّ} \textsuperscript{(69)}،
وهي تفغر \VUgras{مناقيرها} \textsuperscript{(70)} وترفع صياحًا قلقًا.

%%K t07-c2-14-2 p
أما \VUgras{الديك} \textsuperscript{(71)}، وهو أكثر حربية، فيرفع \VUgras{عرفه}
\textsuperscript{(72)} الأحمر، وينفض ريش \VUgras{ذيله} \textsuperscript{(73)}، ويطلق
صياحًا رنّانًا.

%%K t07-c2-15-1 p
15. --- و\VUgras{دجاجة أم} \textsuperscript{(74)} تنقر في \VUgras{كومة الزبل}
\textsuperscript{(81)} مع \VUgras{كتاكيتها} \textsuperscript{(75)}. و\VUgras{الخنزير
الصغير} \textsuperscript{(78)} يجري إلى أمه، \VUgras{الخنزيرة} \textsuperscript{(77)}
الضخمة التي نراها قرب الزريبة.

%%K t07-c2-16-1 p
16. --- وفي أقفاصها تأكل \VUgras{الأرانب} \textsuperscript{(79)} بنهم
\VUgras{العشب} \textsuperscript{(80)} الغضّ الذي تجلبه لها ابنة صاحب المزرعة.
وجِني تحب أرانبها كثيرًا، وفي كل يوم، بعد أن تأخذ \VUgras{البيض}
\textsuperscript{(76)} الطازج من قنّ الدجاج، تأتي لتزور أصدقاءها الصغار.
\VUsaut{6.0mm}
\VUfilet{22mm}
